\documentclass[aspectratio=169,10pt]{beamer}
\usetheme[numbering=fraction,progressbar=frametitle]{metropolis}
\usepackage{fontspec}
\setsansfont{Latin Modern Sans}
\usepackage{amsmath,amssymb,mathtools,graphicx,booktabs}
\usepackage{polyglossia}
\setdefaultlanguage{french}
\setbeamertemplate{navigation symbols}{}
\setbeamerfont{frametitle}{size=\large}
\setbeamerfont{normal text}{size=\normalsize}
\setlength{\parskip}{.45em}
\hypersetup{pdfauthor={},pdfsubject={Lecture slides},pdfkeywords={},pdftitle={Fonctions numériques}}
\providecommand{\N}{\mathbb{N}}
\providecommand{\Z}{\mathbb{Z}}
\providecommand{\D}{\mathbb{D}}
\providecommand{\Q}{\mathbb{Q}}
\providecommand{\R}{\mathbb{R}}
\providecommand{\C}{\mathbb{C}}
\providecommand{\vect}[1]{\overrightarrow{#1}}
\providecommand{\vi}{\vec{\imath}}
\providecommand{\vj}{\vec{\jmath}}
\providecommand{\vk}{\vec{k}}
\providecommand{\norm}[1]{\left\lVert #1\right\rVert}
\providecommand{\vc}[2]{\begin{pmatrix}#1\\#2\end{pmatrix}}
\providecommand{\deter}[2]{\det\left(#1,#2\right)}
\providecommand{\abs}[1]{\left\lvert #1\right\rvert}
\providecommand{\dd}{\,\mathrm{d}}
\providecommand{\e}{\mathrm{e}}
\providecommand{\Card}{\operatorname{Card}}
\providecommand{\E}{\operatorname{E}}
\providecommand{\V}{\operatorname{V}}
\providecommand{\degre}{^{\circ}}
\providecommand{\un}[1]{\,\mathrm{#1}}
\providecommand{\sci}[2]{#1\times10^{#2}}
\providecommand{\pgcd}{\operatorname{PGCD}}
\providecommand{\ppcm}{\operatorname{PPCM}}
\title{Fonctions numériques}
\subtitle{Cours et exemples}
\date{}
\author{}
\begin{document}
\begin{frame}[plain]
\vfill
{\usebeamerfont{title}\inserttitle\par}
\medskip
{\usebeamerfont{subtitle}\insertsubtitle\par}
\medskip
{\usebeamercolor[fg]{progress bar}\rule{\textwidth}{.5pt}}
\vfill
\end{frame}
\begin{frame}{}
Une machine associe à \(x\) le nombre \(x^2-4\). Construire une table pour sept
valeurs de \(x\), placer les points correspondants, puis répondre aux mêmes
questions à partir de la formule, du tableau et du graphique. Comparer ce que
chaque représentation rend facile ou difficile.
\end{frame}
\begin{frame}{Objectifs}
\begin{itemize}

  \item Lire et interpréter une courbe représentative.
  
\item Déterminer domaine de définition, image, antécédent, extremums.
  
\item Dresser et lire un tableau de variation.
  
\item Connaître les fonctions de référence.
  
\item Déduire les courbes des fonctions associées par transformations.

\end{itemize}
\end{frame}
\begin{frame}{Généralités}
Une \textbf{fonction numérique} \(f\) est une application d'un ensemble \(D_f \subset \mathbb{R}\)
(domaine de définition) vers \(\mathbb{R}\). \(f(x)\) est l'\textbf{image} de \(x\),
et \(x\) est un \textbf{antécédent} de \(f(x)\).
\end{frame}
\begin{frame}{Généralités}
Le domaine \(D_f\) doit être symétrique par rapport à \(0\) :
\(x\in D_f\Rightarrow -x\in D_f\). Les égalités suivantes valent pour tout \(x\in D_f\).
\(f\) est \textbf{paire} si \(f(-x)=f(x)\) (courbe symétrique par rapport à \(Oy\)).

\(f\) est \textbf{impaire} si \(f(-x)=-f(x)\) (courbe symétrique par rapport à \(O\)).
\end{frame}
\begin{frame}{Transformations}
À partir de la courbe \(\mathcal{C}_f\) de \(f\) :

\begin{itemize}

  \item \(f(x)+a\) : translation verticale de vecteur \((0,a)\).
  
\item \(f(x+a)\) : translation horizontale de vecteur \((-a,0)\).
  
\item \(-f(x)\) : symétrie par rapport à \(Ox\).
  
\item \(f(-x)\) : symétrie par rapport à \(Oy\).
  
\item \(|f(x)|\) : replier la partie négative vers le haut.

\end{itemize}
\end{frame}
\begin{frame}{Lecture de courbe et vocabulaire}
Le \textbf{domaine de définition} \(D_f\) est l'ensemble des réels \(x\)
pour lesquels \(f(x)\) est défini. Deux causes d'exclusion :

\begin{itemize}

  \item Division par zéro : exclure les \(x\) qui annulent le dénominateur.
  
\item Racine carrée d'un négatif : exclure les \(x\) où l'expression
        sous le radical est \(<0\).

\end{itemize}
\end{frame}
\begin{frame}{Lecture de courbe et vocabulaire}
Pour \(y = f(x)\) :

\begin{itemize}

  \item \(y\) est l'\textbf{image} de \(x\) par \(f\).
  
\item \(x\) est un \textbf{antécédent} de \(y\) : il peut y en avoir
        plusieurs, un seul, ou aucun.

\end{itemize}

\emph{Graphiquement :} l'image de \(x=a\) est l'ordonnée de la courbe
à l'abscisse \(a\) ; les antécédents de \(y=b\) sont les abscisses
des points d'intersection de la courbe avec la droite horizontale \(y=b\).
\end{frame}
\begin{frame}{Lecture de courbe et vocabulaire}
\begin{itemize}
\item \(f\) est \textbf{croissante} sur \(I\) si :
        \(\forall x_1,x_2\in I,\ x_1<x_2\Rightarrow f(x_1)\leq f(x_2)\).
  
\item \(f\) est \textbf{décroissante} sur \(I\) si :
        \(\forall x_1,x_2\in I,\ x_1<x_2\Rightarrow f(x_1)\geq f(x_2)\).
  
\item On dit \textbf{strictement} croissante ou décroissante si les
        inégalités entre les images sont strictes.

\end{itemize}
\end{frame}
\begin{frame}{Lecture de courbe et vocabulaire (2)}
\begin{itemize}
\item \(f(a)\) est un minimum sur \(I\) si \(a\in I\) et
        \(f(a)\leq f(x)\) pour tout \(x\in I\); pour un maximum, \(f(a)\geq f(x)\).
        Un extremum seulement comparé aux valeurs voisines est dit \textbf{local}.

\end{itemize}
\end{frame}
\begin{frame}{Lecture de courbe et vocabulaire}
\begin{enumerate}
\setcounter{enumi}{0}
  \item Identifier le domaine de définition.
  
\item Déterminer les intervalles de croissance et décroissance
        (graphiquement, ou en utilisant le signe de \(f(x_2)-f(x_1)\)).
  
\item Indiquer les valeurs aux bornes et aux extremums.

\end{enumerate}
\end{frame}
\begin{frame}{Lecture de courbe et vocabulaire}
\begin{center}\begin{tabular}{llllll}
\(x\) & \(-\infty\) &  & \(0\) &  & \(+\infty\) \\ 
\(f(x)\) & \(+\infty\) &  & \(-4\) &  & \(+\infty\) \\ 
\end{tabular}\end{center}

\(f\) est décroissante sur \(]-\infty,0]\) et croissante sur \([0,+\infty[\) ;
minimum \(f(0)=-4\).
\end{frame}
\begin{frame}{Transformations de courbes}
À partir de la courbe \(\mathcal{C}_f\) de \(f\) :
\end{frame}
\begin{frame}{Transformations de courbes (2)}
\begin{center}\begin{tabular}{lll}
Nouvelle fonction & Transformation & Description \\ 
\(f(x)+k\) & Translation verticale & Monte de \(k\) (ou descend si \(k<0\)) \\ 
\(f(x+k)\) & Translation horizontale & Glisse à gauche de \(k\) (si \(k>0\)) \\ 
\(-f(x)\) & Symétrie par rapport à \(Ox\) & Retourne la courbe \\ 
\(f(-x)\) & Symétrie par rapport à \(Oy\) & Image miroir \\ 
\(|f(x)|\) & Valeur absolue & Remonte la partie négative \\ 
\end{tabular}\end{center}
\end{frame}
\begin{frame}{Transformations de courbes}
\begin{itemize}

  \item \(x^2+3\) : parabole montée de \(3\) (sommet en \((0,3)\)).
  
\item \((x-2)^2\) : parabole déplacée à droite de \(2\) (sommet en \((2,0)\)).
  
\item \(-(x+1)^2+2\) : parabole retournée, sommet en \((-1,2)\), maximum \(2\).

\end{itemize}
\end{frame}
\end{document}
